% Part: model-theory
% Chapter: models-of-arithmetic
% Section: non-standard-models

\documentclass[../../../include/open-logic-section]{subfiles}

\begin{document}

\olfileid[hi]{mod}{mar}{mdq}
\section{$\Th{Q}$ के प्रतिरूप}

\begin{explain}
हम जानते हैं कि ऐसी अमानक !!{structure}s हैं जो उन्हीं !!{sentence}s को सत्य
बनाती हैं जिन्हें~$\Struct{N}$ सत्य बनाती है, अर्थात् वे~$\Th{TA}$ की
प्रतिरूप हैं। चूँकि $\Sat{N}{\Th{Q}}$,~$\Th{TA}$ का प्रत्येक प्रतिरूप
~$\Th{Q}$ का भी प्रतिरूप है। $\Th{Q}$,~$\Th{TA}$ से बहुत दुर्बल है;
उदाहरणतः
$\Th{Q} \Proves/ \lforall[x][\lforall[y][\eq[(x +
      y)][(y+x)]]]$। दुर्बल सिद्धांतों को संतुष्ट करना अधिक सरल होता है:
उनके अधिक प्रतिरूप होते हैं। उदाहरणतः, $\Th{Q}$ के ऐसे प्रतिरूप हैं जो
$\lforall[x][\lforall[y][\eq[(x + y)][(y+x)]]]$ को असत्य बनाते हैं, किंतु
वे~$\Th{TA}$ के प्रतिरूप नहीं हो सकते, और न ही~$\Th{PA}$ के। $\Th{Q}$ के
प्रतिरूप अपेक्षाकृत सरल भी हैं: हम उन्हें स्पष्ट रूप से निर्दिष्ट कर सकते हैं।
\end{explain}

\begin{ex}
\ollabel{ex:model-K-of-Q}
उस !!{structure}~$\Struct{K}$ पर विचार करें जिसका प्रांत
$\Domain{K} = \Nat
\cup \{a\}$ और जिसकी व्याख्याएँ निम्न हैं:
\begin{align*}
  \Assign{\Obj{0}}{K} & = 0\\
  \Assign{\prime}{K}(x) & =
  \begin{cases}
    x+1 & \text{यदि $x\in \Nat$}\\
    a & \text{यदि $x = a$}
  \end{cases}\\
  \Assign{+}{K}(x, y) & =
  \begin{cases}
    x+y & \text{यदि $x$, $y \in\Nat$}\\
    a & \text{अन्यथा}
  \end{cases}\\
  \Assign{\times}{K}(x, y) & =
  \begin{cases}
    xy & \text{यदि $x$, $y \in\Nat$}\\
    0 & \text{यदि $x = 0$ या $y = 0$}\\
    a & \text{अन्यथा}\\
  \end{cases}\\
  \Assign{<}{K} & =
  \Setabs{\tuple{x,y}}{x, y \in \Nat \text{ और } x<y} \cup
  \Setabs{\tuple{x,a}}{x \in \Domain{K}}
\end{align*}
$\Sat{K}{\Th{Q}}$ दिखाने के लिए हमें सत्यापित करना है कि~$\Th{Q}$ के सभी
अभिगृहीत~$\Struct{K}$ में सत्य हैं। सुविधा के लिए, $x^\nssucc$ को
$\Assign{\prime}{K}(x)$ के लिए (~$x$ का,~$\Struct{K}$ में, “परवर्ती”),
$x \nsplus y$ को $\Assign{+}{K}(x, y)$ के लिए (~$x$ और $y$ का,
~$\Struct{K}$ में, “योग”), $x \nstimes y$ को
$\Assign{\times}{K}(x, y)$ के लिए (~$x$ और~$y$ का,~$\Struct{K}$ में,
“गुणनफल”), और $x \nsless y$ को
$\tuple{x,y} \in
\Assign{<}{K}$ के लिए लिखें। इन संक्षेपों से
हम~$\Struct{K}$ की संक्रियाएँ अधिक स्पष्ट रूप में इस प्रकार दे सकते हैं:
\[
\begin{array}{c|c}
  x & x^\nssucc \\
  \hline
  n & n+1 \\
  a & a
\end{array}
\qquad
\begin{array}{c|ccc}
  x \nsplus y & 0 & m & a \\
  \hline
  0 & 0 & m & a \\
  n & n & n+m & a \\
  a & a & a & a \\
\end{array}
\qquad
\begin{array}{c|ccc}
  x \nstimes y & 0 & m & a \\
  \hline
  0 & 0 & 0 & 0 \\
  n & 0 & nm & a \\
  a & 0 & a & a \\
\end{array}
\]
हमारे पास $n \nsless m$ तभी और केवल तभी जब $n<m$, जहाँ $n$,
$m \in \Nat$; और $x \nsless
a$ प्रत्येक~$x \in \Domain{K}$ के लिए।

$\Sat{K}{\lforall[x][\lforall[y][(\eq[x'][y'] \lif \eq[x][y])]]}$,
क्योंकि $\nssucc$ !!{injective} है।
$\Sat{K}{\lforall[x][\eq/[\Obj
0][x']]}$, क्योंकि $0$ कोई $\nssucc$-परवर्ती नहीं है~$\Struct{K}$ में।
$\Sat{K}{\lforall[x][(\eq[x][\Obj 0] \lor \lexists[y][\eq[x][y']])]}$,
क्योंकि प्रत्येक $n>0$ के लिए $n = (n-1)^\nssucc$, और
$a = a^\nssucc$।

$\Sat{K}{\lforall[x][\eq[(x + \Obj 0)][x]]}$, क्योंकि
$n \nsplus 0 = n+0
= n$, और $a\nsplus 0 = a$,~$\nsplus$ की परिभाषा से।
$\Sat{K}{\lforall[x][\lforall[y][\eq[(x + y')][(x + y)']]]}$ थोड़ा अधिक
जटिल है। यदि $n$, $m$ दोनों मानक हों, तो हमारे पास:
\begin{align*}
(n \nsplus m^\nssucc) & = (n+(m+1)) = (n+m)+1 = (n \nsplus m)^\nssucc  
\intertext{क्योंकि मानक संख्याओं पर $\nsplus$ और $^\nssucc$ क्रमशः $+$
  और $\prime$ से मेल खाते हैं। अब मान लें कि $x \in \Domain{K}$। तब}
(x \nsplus a^\nssucc) & = (x \nsplus a) = a = a^\nssucc = (x \nsplus a)^\nssucc
\intertext{शेष स्थिति तब है जब $y \in \Domain{K}$, किंतु $x =
  a$। यहाँ हमें इस आधार पर स्थितियाँ अलग करनी होंगी कि $y =
  n$ मानक है या $y = b$:}
(a \nsplus n^\nssucc) & = (a \nsplus (n+1)) = a = a^\nssucc = (a \nsplus n)^\nssucc\\
(a \nsplus a^\nssucc) & = (a \nsplus a) = a = a^\nssucc = (a \nsplus a)^\nssucc
\end{align*}
यह निस्संदेह आवश्यकता से थोड़ा अधिक विस्तृत है। उदाहरणार्थ, चूँकि
$a \nsplus z = a$,~$z$ चाहे जो हो, इसलिए हम तुरंत निष्कर्ष निकाल सकते हैं
कि $a \nsplus
a^\nssucc = a$। शेष अभिगृहीतों को भी इसी प्रकार सत्यापित
किया जा सकता है।

इस प्रकार $\Struct{K}$,~$\Th{Q}$ का प्रतिरूप है। उसका “योग”~$\nsplus$
क्रमविनिमेय भी है। किंतु कुछ अन्य वाक्य~$\Struct{N}$ में सत्य और
~$\Struct{K}$ में असत्य हैं, तथा इसके विपरीत भी। उदाहरणार्थ,
$a \nsless a$, अतः $\Sat{K}{\lexists[x][x < x]}$ और
$\Sat/{K}{\lforall[x][\lnot x<x]}$। इससे दिखता है कि
$\Th{Q} \Proves/
\lforall[x][\lnot x < x]$।
\end{ex}

\begin{prob}
सिद्ध कीजिए कि \olref[mod][mar][mdq]{ex:model-K-of-Q} की
$\Struct{K}$,~$\Th{Q}$ के शेष अभिगृहीतों को संतुष्ट करती है:
\begin{align*}
  & \lforall[x][\eq[(x \times \Obj 0)][\Obj 0]] \tag{$!Q_6$}\\
  & \lforall[x][\lforall[y][\eq[(x \times y')][((x \times y) + x)]]] \tag{$!Q_7$}\\
  & \lforall[x][\lforall[y][(x < y \liff \lexists[z][\eq[(z' + x)][y]])]] \tag{$!Q_8$}
\end{align*}
केवल~$\prime$ को अंतर्विष्ट करने वाला कोई ऐसा !!a{sentence} खोजिए जो
~$\Struct{N}$ में सत्य, किंतु~$\Struct{K}$ में असत्य हो।
\end{prob}

\begin{ex}
\ollabel{ex:model-L-of-Q} उस !!{structure}~$\Struct{L}$ पर विचार करें जिसका
प्रांत $\Domain{L} = \Nat \cup \{a, b\}$ है और जिसकी व्याख्याएँ
$\Assign{\prime}{L} = \nssucc$, $\Assign{+}{L} = \nsplus$ निम्न से दी
गई हैं:
\[
\begin{array}{c|c}
  x & x^\nssucc \\
  \hline
  n & n+1 \\
  a & a\\
  b & b
\end{array}
\qquad
\begin{array}{c|ccc}
  x \nsplus y & m & a & b\\
  \hline
  n & n+m & b & a\\
  a & a & b & a\\
  b & b & b & a
\end{array}
\]
चूँकि $\nssucc$ !!{injective} है, $0$ उसके परास में नहीं है, और प्रत्येक
$x \in \Domain{L}$,~$0$ के अतिरिक्त, उसके परास में है, इसलिए
अभिगृहीत $!Q_1$--$!Q_3$,~$\Struct{L}$ में सत्य हैं। प्रत्येक $x$ के लिए
$x \nsplus 0 = x$, इसलिए $!Q_4$ भी सत्य है। $!Q_5$ के लिए
$x \nsplus y^\nssucc$ और $(x \nsplus
y)^\nssucc$ पर विचार करें। यदि
$x$ और $y$ दोनों मानक हों, तो वे बराबर हैं, क्योंकि तब $\nssucc$ और
$\nsplus$ क्रमशः $\prime$ और $+$ से मेल खाते हैं। यदि $x$ अमानक और $y$
मानक हो, तो $x \nsplus y^\nssucc = x =
x^\nssucc = (x \nsplus y)^\nssucc$। यदि $x$ और $y$ दोनों अमानक हों, तो
चार स्थितियाँ हैं:
\begin{align*}
&  a \nsplus a^\nssucc  = b = b^\nssucc = (a \nsplus a)^\nssucc\\
&  b \nsplus b^\nssucc  = a = a^\nssucc = (b \nsplus b)^\nssucc\\
&  b \nsplus a^\nssucc  = b = b^\nssucc = (b \nsplus y)^\nssucc\\
&  a \nsplus b^\nssucc  = a = a^\nssucc = (a \nsplus b)^\nssucc\\
\intertext{यदि $x$ मानक, किंतु $y$ अमानक हो, तो हमारे पास}
&  n \nsplus a^\nssucc  = n \nsplus a = b = b^\nssucc = (n \nsplus a)^\nssucc\\
&  n \nsplus b^\nssucc  = n \nsplus b = a = a^\nssucc = (n \nsplus b)^\nssucc
\end{align*}
अतः $\Sat{L}{!Q_5}$। किंतु $a \nsplus 0 \neq 0 \nsplus a$, इसलिए
$\Sat/{L}{\lforall[x][\lforall[y][\eq[(x+y)][(y+x)]]]}$।
\end{ex}

\begin{prob}
\olref[mod][mar][mdq]{ex:model-L-of-Q} की $\Struct{L}$ को ऐसे
$\nstimes$ और $\nsless$ सम्मिलित करके विस्तृत कीजिए जो~$\times$ और $<$ की
व्याख्या करें। दिखाइए कि आपकी संरचना~$\Th{Q}$ के शेष अभिगृहीतों को संतुष्ट
करती है:
\begin{align*}
& \lforall[x][\eq[(x \times \Obj 0)][\Obj 0]] \tag{$!Q_6$}\\ &
  \lforall[x][\lforall[y][\eq[(x \times y')][((x \times y) + x)]]]
  \tag{$!Q_7$}\\ & \lforall[x][\lforall[y][(x < y \liff \lexists[z][\eq[(z'+x)][y]])]] \tag{$!Q_8$}
\end{align*}
\end{prob}

\begin{prob}
\olref[mod][mar][mdq]{ex:model-L-of-Q} की $\Struct{L}$ में
$a^\nssucc
= a$ और $b^\nssucc = b$। क्या~$\Th{Q}$ का कोई ऐसा प्रतिरूप
है जिसमें $a^\nssucc = b$ और $b^\nssucc = a$?
\end{prob}

\begin{explain}
हमने~$\Th{Q}$ के ऐसे प्रतिरूप स्पष्ट रूप से बनाए हैं जिनमें अमानक
!!{element}s मानक अवयवों के “आगे” रहते हैं। वस्तुतः, अभिगृहीतों से इतना होना
आवश्यक है। कोई अमानक !!{element}~$x$, ${} \nsless 0$ नहीं हो सकता, क्योंकि
$\Th{Q} \Proves
\lforall[x][\lnot x<0]$
(\olref[inc][req][min]{lem:less-zero} देखें)। साथ ही प्रत्येक $n$ के लिए,
$\Th{Q} \Proves \lforall[x][(x < \num{n}' \lif
(\eq[x][\num{0}] \lor \eq[x][\num{1}] \lor \dots \lor
\eq[x][\num{n}]))]$ (\olref[inc][req][min]{lem:less-nsucc}), इसलिए
$a \nsless n$ किसी भी~$n>0$ के लिए नहीं हो सकता।
\end{explain}

\end{document}
